
\section{Ie\c{s}an--Saint Venant Formulation}

\subsection{Saint-Venant class (relaxed elastostatics)}

We consider a prismatic, homogeneous, and isotropic linearly elastic cylinder
\(
\Body=\Section\times(0,L)
\),
where \(\Section\subset\mathbb{R}^2\) is a bounded (possibly multiply connected) cross-section
spanned by \(\{\FrameBasis_y,\FrameBasis_z\}\) and the axis is \(\FrameBasis\).
Set \(\bm{r}=y\,\FrameBasis_y+z\,\FrameBasis_z\) and
\(
\mathbf{1}=\FrameBasis\otimes\FrameBasis+\mathbf{1}_\Section
\)
with
\(
\mathbf{1}_\Section=\FrameBasis_y\otimes\FrameBasis_y+\FrameBasis_z\otimes\FrameBasis_z
\).
The lateral boundary is \(\BodyLateralSurface=\SectionBoundary\times(0,L)\).

\begin{remark}
The developments are frame-free (only right-handedness is assumed). For implementation it is convenient
either to take the first coordinate along \(\FrameBasis\) (so \(\FrameBasis\equiv\bm\partial_x\)) or the last;
both conventions are compatible with the formulas below.
\end{remark}


The elastostatic field \((\bm u,\bm T)\) satisfies
\begin{equation}
\left\{
\begin{aligned}
\Div\,\bm T&=\bm 0,\\
\bm T&=\lambda\,(\operatorname{tr}\bm\varepsilon)\,\mathbf{1}+2G\,\bm\varepsilon,\\
\bm\varepsilon&=\tfrac12\big(\nabla\bm u+\nabla\bm u^{\!\top}\big),
\end{aligned}\right.
\qquad\text{in }\Body,
\label{eq:3D-LE}
\end{equation}
together with \(\bm T\,\bm n=\bm 0\) on \(\BodyLateralSurface\).
In the \emph{relaxed} Saint-Venant problem, end tractions are replaced by resultants
\(
\mathscr{F}(\bm u)\in\mathbb{R}^3
\)
and
\(
\mathscr{M}(\bm u)\in\mathbb{R}^3
\) evaluated on, say, \(\Section_1=\Section\times\{x=0\}\):
\begin{equation}
\mathscr{F}(\bm u)=\int_{\Section_1}\bm T(\bm u)\,(-\FrameBasis)\mathrm{~d}A,\qquad
\mathscr{M}(\bm u)=\int_{\Section_1}\bm{r}\times\big(\bm T(\bm u)\,(-\FrameBasis)\big)\mathrm{~d}A.
\label{eq:def-resultants}
\end{equation}
We adopt the standard Saint-Venant stress representation
\begin{equation}
\boxed{%
\bm T=\AxialStress\,\FrameBasis\otimes\FrameBasis
+\FrameBasis\otimes\bm{\tau}+\bm{\tau}\otimes\FrameBasis,
}
\qquad
\AxialStress:=\FrameBasis\cdot\bm T\,\FrameBasis,\quad
\bm{\tau}:=(\PlaneBasis_\beta\otimes\PlaneBasis_\beta)\,\bm T\,\FrameBasis,
\label{eq:saint-venant-stress}
\end{equation}
and decompose the relaxed class into:
\[
(P_1):\ \text{extension--bending--torsion with }F_\alpha=0,\qquad
(P_2):\ \text{flexure with }F_3=M_i=0\ \text{on }\Section_1.
\]

\subsection{Problem \(P_1\): extension--bending--torsion}
\subsubsection*{Characterization}
Besides \eqref{eq:3D-LE} and \(\bm T\,\bm n=\bm 0\) on \(\BodyLateralSurface\), we impose
\[
F_\alpha=0\quad(\alpha=1,2),\qquad
\bm u'\ \text{is rigid on each section, i.e.}\ \ \bm u'=\bm\alpha+\bm\beta\times\bm{r},
\]
with constants \(a_\varepsilon=\FrameBasis\cdot\bm\alpha\),
\(a_\vartheta=\FrameBasis\cdot\bm\beta\),
\(\bm a_\Section=\FrameBasis\times\bm\beta_\kappa\).
This yields the section warping \(\WarpTwistIesan\) (torsion) by the Neumann problem
\begin{equation}
\boxed{%
\Delta\WarpTwistIesan=0\ \text{ in }\Section,\qquad
\partial_{\bm\nu}\WarpTwistIesan=\varepsilon_{\alpha\beta}n_\alpha x_\beta\ \text{ on }\SectionBoundary,
}
\label{eq:bvp-warp-twist}
\end{equation}
(unique up to an additive constant).

\subsubsection*{Sectional displacement, strain, stress}
The displacement induced by \(\bm a=(\bm a_\Section,a_\varepsilon,a_\vartheta)\) is
\begin{equation}
\boxed{%
\begin{aligned}
\hat{\bm u}\{\bm a\}
&=\Big(-\tfrac12 x^2\,\mathbf{1}-\nu\,\operatorname{dev}\bm{r}\otimes\bm{r}+x\,\FrameBasis\otimes\bm{r}\Big)\bm a_\Section\\
&\quad+(x\,\FrameBasis-\nu\,\bm{r})\,a_\varepsilon
+\big(x\,\FrameBasis\times\bm{r}+\WarpTwistIesan\,\FrameBasis\big)a_\vartheta .
\end{aligned}}
\label{eq:iesan-displacement-1}
\end{equation}
The corresponding strain decomposes as
\begin{equation}
\boxed{%
\bm E^{(a)}=
\big(a_\varepsilon+\bm{r}\cdot\bm a_\Section\big)\,(\FrameBasis\otimes\FrameBasis-\nu\,\mathbf{1}_\Section)
+\tfrac12\big(\FrameBasis\otimes\bm\varepsilon_\vartheta+\bm\varepsilon_\vartheta\otimes\FrameBasis\big),}
\quad
\bm\varepsilon_\vartheta:=a_\vartheta\big(\FrameBasis\times\bm{r}+\nabla\WarpTwistIesan\big),
\label{eq:iesan-strain-tensor-1}
\end{equation}
and the axial/shear stresses take the Saint-Venant form
\begin{equation}
\boxed{%
\bm{\tau}=G\big(\nabla\WarpTwistIesan-\FrameBasis\times\bm{r}\big)a_\vartheta,\qquad
\AxialStress=E\big(\bm{r}\cdot\bm a_\Section+a_\varepsilon\big).}
\label{eq:iesan-stress-1}
\end{equation}

\subsubsection*{Resultants and torsional constant}
On \(\Section_1\) with outward normal \(-\FrameBasis\) the resultants give the linear system
\begin{equation}
\boxed{%
E\big(\IntRoR\,\bm a_\Section+A\,a_\varepsilon\,\CentroidLocation\big)=\FrameBasis\times\bm M,\qquad
EA\big(\bm a_\Section\cdot\CentroidLocation+a_\varepsilon\big)=-N,\qquad
G\,\Jo\,a_\vartheta=-T,}
\label{eq:2-15}
\end{equation}
with
\(
\IntRoR=\int_\Section \bm{r}\otimes\bm{r}\mathrm{~d}A
\),
area \(A\), centroid \(\CentroidLocation\),
and torsion constant
\begin{equation}
\boxed{%
\Jo=\int_\Section\Big(r^2+(\FrameBasis\times\bm{r})\cdot\nabla\WarpTwistIesan\Big)\mathrm{~d}A.}
\label{eq:def-j-origin}
\end{equation}

\subsection{Problem \(P_2\): flexure}
\subsubsection*{Reduction and data}
If \(\bm u\) solves \(P_2\) with shear resultant \(\ShearForce\) on \(\Section_1\),
then \(\bm u'\) solves \(P_1\) with bending moment \(\bm M=\FrameBasis\times\ShearForce\).
Consequently the parameters \(\bm b=(\bm b_\Section,b_\varepsilon,b_\vartheta)\) entering the flexure field satisfy
\begin{equation}
\boxed{%
E\big(\IntRoR\,\bm b_\Section+A\,b_\varepsilon\,\CentroidLocation\big)=-\ShearForce,\qquad
\bm b_\Section\cdot\CentroidLocation+b_\varepsilon=0,\qquad
b_\vartheta=0,}
\label{eq:2-18}
\end{equation}
i.e.
\(
\bm b_\Section=-E^{-1}\big(\IntRoR-A\,\CentroidLocation\!\otimes\!\CentroidLocation\big)^{-1}\ShearForce
\)
and \(b_\varepsilon=-\CentroidLocation\cdot\bm b_\Section\).

\subsubsection*{Flexural warping}
The (scalar) flexure warping \(\WarpShearIesan\) is determined by the Poisson--Neumann problem
\begin{equation}
\boxed{%
\left\{\begin{aligned}
\Delta\WarpShearIesan&=-2\big((\bm{r}-\CentroidLocation)\cdot\bm b_\Section\big)\quad &&\text{in }\Section,\\
\partial_{\bm\nu}\WarpShearIesan&=\nu\,\bm\nu\cdot\Big(\operatorname{dev}\bm{r}\otimes\bm{r}-\bm{r}\otimes\CentroidLocation\Big)\bm b_\Section
\quad &&\text{on }\SectionBoundary,\\
\int_\Section \WarpShearIesan\mathrm{~d}A&=0, &&
\end{aligned}\right.}
\label{eq:bvp-warp-shear-iesan}
\end{equation}
whose compatibility is equivalent to \(\bm b_\Section\cdot\CentroidLocation+b_\varepsilon=0\), fixing \(\CentroidLocation\) at the area centroid.

\subsubsection*{Sectional displacement, strain, stress}
Up to rigid motion, the flexure displacement reads
\begin{equation}
\boxed{%
\begin{aligned}
\hat{\bm u}_\Section(x,\bm{r})
&=-\tfrac16\,x^3\,\bm b_\Section
+x\,\nu\Big(\tfrac12 r^2\,\bm b_\Section-(\bm b_\Section\cdot\bm{r})\,\bm{r}\Big)
+x\,\nu\,(\CentroidLocation\cdot\bm b_\Section)\,\bm{r}
+c_\vartheta\,x\,(\FrameBasis\times\bm{r}),\\
\FrameBasis\cdot\hat{\bm u}(x,\bm{r})
&=\tfrac12\,x^2\,\big((\bm{r}-\CentroidLocation)\cdot\bm b_\Section\big)+c_\vartheta\,\WarpTwistIesan(\bm{r})+\WarpShearIesan(\bm{r}),
\end{aligned}}
\label{eq:iesan-displacement-2}
\end{equation}
with \(c_\vartheta\) fixed by
\begin{equation}
\boxed{%
\Jo\,c_\vartheta
=-\int_\Section\Big[(\FrameBasis\times\bm{r})\cdot\nabla\big(\bm{\WarpShearIesan}\cdot\bm b_\Section\big)
+\tfrac12\,\nu\,r^2\,(\FrameBasis\times\bm{r})\cdot\bm b_\Section\Big]\mathrm{d}A.}
\label{eq:iesan-2-22}
\end{equation}
The associated infinitesimal strain admits the compact form
\begin{equation}
\boxed{%
\begin{aligned}
\bm E(\hat{\bm u})
&=(\varepsilon\,\FrameBasis+\bm\varepsilon_\gamma)\stackrel{\rm s}{\otimes}\FrameBasis
-\nu\,\varepsilon\,\mathbf{1}_\Section,\\
\varepsilon&=x\,(\bm{r}-\CentroidLocation)\cdot\bm b_\Section,\\
\bm\varepsilon_\gamma
&=(a_\vartheta+c_\vartheta)\big(\nabla\WarpTwistIesan+\FrameBasis\times\bm{r}\big)
-\nu\Big(\operatorname{dev}(\bm{r}\otimes\bm{r})-\bm{r}\otimes\CentroidLocation\Big)\bm b_\Section
+\nabla\big(\bm{\WarpShearIesan}\cdot\bm b_\Section\big),
\end{aligned}}
\label{eq:iesan-strain-2}
\end{equation}
and the Saint-Venant axial/shear stress components are
\begin{equation}
\boxed{%
\AxialStress(\hat{\bm u})=E\,x\,(\bm{r}-\CentroidLocation)\cdot\bm b_\Section,\qquad
\bm{\tau}(\hat{\bm u})
=G\Big[(a_\vartheta+c_\vartheta)\big(\nabla\WarpTwistIesan+\FrameBasis\times\bm{r}\big)
-\nu\Big(\operatorname{dev}(\bm{r}\otimes\bm{r})-\bm{r}\otimes\CentroidLocation\Big)\bm b_\Section
+\nabla\big(\bm{\WarpShearIesan}\cdot\bm b_\Section\big)\Big].}
\label{eq:iesan-stress-2}
\end{equation}

\subsection{Equilibrium identities in the Saint-Venant class}
For any superposition satisfying \eqref{eq:saint-venant-stress} and \(\bm T\,\bm n=\bm 0\) on \(\BodyLateralSurface\),
the field equations reduce on each section to
\begin{subequations}
\label{eq:sv-equilibrium}
\begin{align}
\bm{\tau}'&=\bm 0 && \text{in }\Section, \label{eq:shear_x}\\
\DivS\bm{\tau}+\AxialStress'&=0 && \text{in }\Section, \\
\bm{\tau}\cdot\bm\nu&=0 && \text{on }\SectionBoundary.
\end{align}
\end{subequations}
Thus \(\bm{\tau}=\bm{\tau}(\bm{r})\) is \(x\)-independent by \eqref{eq:shear_x}, and, for flexure, \(\DivS\bm{\tau}=-E\,(\bm{r}-\CentroidLocation)\cdot\bm b_\Section\).

\subsection{General Saint-Venant solution (displacement)}
Collecting \eqref{eq:iesan-displacement-1} and \eqref{eq:iesan-displacement-2}, the full family
parameterized by \(\bm a=(\bm a_\Section,a_\varepsilon,a_\vartheta)\) and \(\bm b_\Section\) reads
\begin{equation}
\boxed{%
\begin{aligned}
\hat{\bm u}\{\bm a,\bm b_\Section\}
&=\Big(-\tfrac12 x^2\,\mathbf{1}-\nu\,\operatorname{dev}\bm{r}\otimes\bm{r}+x\,\FrameBasis\otimes\bm{r}\Big)\bm a_\Section
+(x\,\FrameBasis-\nu\,\bm{r})\,a_\varepsilon
+\big(x\,\FrameBasis\times\bm{r}+\WarpTwistIesan\,\FrameBasis\big)a_\vartheta\\
&\quad-\tfrac16\,x^3\,\bm b_\Section
-x\,\nu\Big(\operatorname{dev}\bm{r}\otimes\bm{r}-\bm{r}\otimes\CentroidLocation\Big)\bm b_\Section
+\tfrac12 x^2\,\FrameBasis\otimes(\bm{r}-\CentroidLocation)\,\bm b_\Section
+\FrameBasis\otimes\bm{\WarpShearIesan}(\bm{r})\,\bm b_\Section\\
&\quad+\big(x\,\FrameBasis\times\bm{r}+\WarpTwistIesan\,\FrameBasis\big)c_\vartheta,
\end{aligned}}
\label{eq:iesan-general-solution}
\end{equation}
with \(\bm a\) fixed by \eqref{eq:2-15}, \(\bm b\) by \eqref{eq:2-18}, \(\WarpTwistIesan\) by \eqref{eq:bvp-warp-twist},
\(\WarpShearIesan\) by \eqref{eq:bvp-warp-shear-iesan}, and \(c_\vartheta\) by \eqref{eq:iesan-2-22}.
Choosing principal centroidal axes reduces the parameter relations to the classical one-dimensional forms.
